\chapter{Boquerón}\label{dept:Boquerón}
\mapaparaguai{16}{Boquerón}
\vfill\clearpage
\section{Boqueron}

\begin{itemize}
\tightlist
\item
  Filadelfia e a capital de facto, polo urbano e industrial do Chaco
  Central mennonita; 156.425 hab. no distrito.
\item
  Vocacao de pecuaria intensiva, laticinios (cooperativas Fernheim,
  Neuland, Menno) e base aerea estratégica em Mariscal Estigarribia.
\item
  Lacuna de solar/\allowbreak pluviometria da capital fechada com dados NASA POWER
  (2001-2020).
\end{itemize}

\subsection{Ficha climática de Filadelfia (capital de facto
departamental)}

\textbf{Irradiância solar global --- (kWh/\allowbreak m²/\allowbreak dia):}

{\def\LTcaptype{none} % do not increment counter
\begin{longtable}[]{@{}
  >{\raggedright\arraybackslash}p{(\linewidth - 22\tabcolsep) * \real{0.0833}}
  >{\raggedright\arraybackslash}p{(\linewidth - 22\tabcolsep) * \real{0.0833}}
  >{\raggedright\arraybackslash}p{(\linewidth - 22\tabcolsep) * \real{0.0833}}
  >{\raggedright\arraybackslash}p{(\linewidth - 22\tabcolsep) * \real{0.0833}}
  >{\raggedright\arraybackslash}p{(\linewidth - 22\tabcolsep) * \real{0.0833}}
  >{\raggedright\arraybackslash}p{(\linewidth - 22\tabcolsep) * \real{0.0833}}
  >{\raggedright\arraybackslash}p{(\linewidth - 22\tabcolsep) * \real{0.0833}}
  >{\raggedright\arraybackslash}p{(\linewidth - 22\tabcolsep) * \real{0.0833}}
  >{\raggedright\arraybackslash}p{(\linewidth - 22\tabcolsep) * \real{0.0833}}
  >{\raggedright\arraybackslash}p{(\linewidth - 22\tabcolsep) * \real{0.0833}}
  >{\raggedright\arraybackslash}p{(\linewidth - 22\tabcolsep) * \real{0.0833}}
  >{\raggedright\arraybackslash}p{(\linewidth - 22\tabcolsep) * \real{0.0833}}@{}}
\toprule\noalign{}
\begin{minipage}[b]{\linewidth}\raggedright
Jan
\end{minipage} & \begin{minipage}[b]{\linewidth}\raggedright
Fev
\end{minipage} & \begin{minipage}[b]{\linewidth}\raggedright
Mar
\end{minipage} & \begin{minipage}[b]{\linewidth}\raggedright
Abr
\end{minipage} & \begin{minipage}[b]{\linewidth}\raggedright
Mai
\end{minipage} & \begin{minipage}[b]{\linewidth}\raggedright
Jun
\end{minipage} & \begin{minipage}[b]{\linewidth}\raggedright
Jul
\end{minipage} & \begin{minipage}[b]{\linewidth}\raggedright
Ago
\end{minipage} & \begin{minipage}[b]{\linewidth}\raggedright
Set
\end{minipage} & \begin{minipage}[b]{\linewidth}\raggedright
Out
\end{minipage} & \begin{minipage}[b]{\linewidth}\raggedright
Nov
\end{minipage} & \begin{minipage}[b]{\linewidth}\raggedright
Dez
\end{minipage} \\
\midrule\noalign{}
\endhead
\bottomrule\noalign{}
\endlastfoot
6.59 & 6.02 & 5.41 & 4.37 & 3.38 & 2.88 & 3.34 & 4.19 & 4.79 & 5.46 &
6.25 & 6.43 \\
\end{longtable}
}

Média anual: 4.92 kWh/\allowbreak m²/\allowbreak dia. Inclinação recomendada: 22° Norte (anual);
32° inverno; 12° verão.

\textbf{Precipitação --- (mm/\allowbreak mês):}

{\def\LTcaptype{none} % do not increment counter
\begin{longtable}[]{@{}
  >{\raggedright\arraybackslash}p{(\linewidth - 22\tabcolsep) * \real{0.0833}}
  >{\raggedright\arraybackslash}p{(\linewidth - 22\tabcolsep) * \real{0.0833}}
  >{\raggedright\arraybackslash}p{(\linewidth - 22\tabcolsep) * \real{0.0833}}
  >{\raggedright\arraybackslash}p{(\linewidth - 22\tabcolsep) * \real{0.0833}}
  >{\raggedright\arraybackslash}p{(\linewidth - 22\tabcolsep) * \real{0.0833}}
  >{\raggedright\arraybackslash}p{(\linewidth - 22\tabcolsep) * \real{0.0833}}
  >{\raggedright\arraybackslash}p{(\linewidth - 22\tabcolsep) * \real{0.0833}}
  >{\raggedright\arraybackslash}p{(\linewidth - 22\tabcolsep) * \real{0.0833}}
  >{\raggedright\arraybackslash}p{(\linewidth - 22\tabcolsep) * \real{0.0833}}
  >{\raggedright\arraybackslash}p{(\linewidth - 22\tabcolsep) * \real{0.0833}}
  >{\raggedright\arraybackslash}p{(\linewidth - 22\tabcolsep) * \real{0.0833}}
  >{\raggedright\arraybackslash}p{(\linewidth - 22\tabcolsep) * \real{0.0833}}@{}}
\toprule\noalign{}
\begin{minipage}[b]{\linewidth}\raggedright
Jan
\end{minipage} & \begin{minipage}[b]{\linewidth}\raggedright
Fev
\end{minipage} & \begin{minipage}[b]{\linewidth}\raggedright
Mar
\end{minipage} & \begin{minipage}[b]{\linewidth}\raggedright
Abr
\end{minipage} & \begin{minipage}[b]{\linewidth}\raggedright
Mai
\end{minipage} & \begin{minipage}[b]{\linewidth}\raggedright
Jun
\end{minipage} & \begin{minipage}[b]{\linewidth}\raggedright
Jul
\end{minipage} & \begin{minipage}[b]{\linewidth}\raggedright
Ago
\end{minipage} & \begin{minipage}[b]{\linewidth}\raggedright
Set
\end{minipage} & \begin{minipage}[b]{\linewidth}\raggedright
Out
\end{minipage} & \begin{minipage}[b]{\linewidth}\raggedright
Nov
\end{minipage} & \begin{minipage}[b]{\linewidth}\raggedright
Dez
\end{minipage} \\
\midrule\noalign{}
\endhead
\bottomrule\noalign{}
\endlastfoot
104 & 127 & 105 & 80 & 63 & 26 & 18 & 12 & 22 & 75 & 106 & 117 \\
\end{longtable}
}

Média anual: \textasciitilde854 mm/\allowbreak ano. Regime semiárido Chaco Central.
Estação chuvosa Nov-Mar; seca Jun-Set.
\clearpage
\subsection*{Dados Departamentais}
\subsubsection{Desenvolvimento Humano e
Social}

\subsection{IDH Departamental}

{\def\LTcaptype{none} % do not increment counter
\begin{longtable}[]{@{}ll@{}}
\toprule\noalign{}
Indicador & Valor \\
\midrule\noalign{}
\endhead
\bottomrule\noalign{}
\endlastfoot
IDH & 0,676 \\
Classificação IDH & médio \\
Ranking nacional (1=melhor) & 16/\allowbreak 18 \\
Esperança de vida & 70,8 anos \\
Anos de escolaridade média & 7,1 anos \\
RNB per capita (USD PPA) & 8.600 \\
\end{longtable}
}

\subsection{Indicadores de Pobreza}

{\def\LTcaptype{none} % do not increment counter
\begin{longtable}[]{@{}ll@{}}
\toprule\noalign{}
Indicador & Valor \\
\midrule\noalign{}
\endhead
\bottomrule\noalign{}
\endlastfoot
Taxa de pobreza total (\%) & estimativa 24,0\% \\
Taxa de pobreza extrema (\%) & estimativa 7,5\% \\
Índice de Gini & estimativa 0,485 \\
\end{longtable}
}

\subsection{Infraestrutura Básica}

{\def\LTcaptype{none} % do not increment counter
\begin{longtable}[]{@{}ll@{}}
\toprule\noalign{}
Indicador & Valor \\
\midrule\noalign{}
\endhead
\bottomrule\noalign{}
\endlastfoot
Acesso a água potável (\%) & estimativa 52,3\% \\
Acesso a saneamento (\%) & estimativa 48,7\% \\
\end{longtable}
}

\subsubsection{Segurança Pública}

\subsection{Indicadores}

{\def\LTcaptype{none} % do not increment counter
\begin{longtable}[]{@{}
  >{\raggedright\arraybackslash}p{(\linewidth - 2\tabcolsep) * \real{0.6111}}
  >{\raggedright\arraybackslash}p{(\linewidth - 2\tabcolsep) * \real{0.3889}}@{}}
\toprule\noalign{}
\begin{minipage}[b]{\linewidth}\raggedright
Indicador
\end{minipage} & \begin{minipage}[b]{\linewidth}\raggedright
Valor
\end{minipage} \\
\midrule\noalign{}
\endhead
\bottomrule\noalign{}
\endlastfoot
Taxa de homicídios (por 100k hab) & 0,7--1,7 (2019--2020, fonte INE\footnote{Instituto Nacional de Estadística (Instituto Nacional de Estatística), órgão oficial de dados demográficos e censitários.} ---
entre as mais baixas do país) \\
Crimes registrados (total/\allowbreak ano) & muito baixo (população esparsa de
\textasciitilde60 mil hab. em vasto território) \\
Número de comisarías & \textasciitilde5 (estimativa --- território
imenso, mas população concentrada em poucas cidades) \\
Índice segurança relativo & alto \\
\end{longtable}
}

\subsubsection{Mercado de Terra Rural}

\subsection{Preços por Tipo de Terra}

{\def\LTcaptype{none} % do not increment counter
\begin{longtable}[]{@{}
  >{\raggedright\arraybackslash}p{(\linewidth - 6\tabcolsep) * \real{0.1538}}
  >{\raggedright\arraybackslash}p{(\linewidth - 6\tabcolsep) * \real{0.3846}}
  >{\raggedright\arraybackslash}p{(\linewidth - 6\tabcolsep) * \real{0.1795}}
  >{\raggedright\arraybackslash}p{(\linewidth - 6\tabcolsep) * \real{0.2821}}@{}}
\toprule\noalign{}
\begin{minipage}[b]{\linewidth}\raggedright
Tipo
\end{minipage} & \begin{minipage}[b]{\linewidth}\raggedright
Preço (USD/\allowbreak ha)
\end{minipage} & \begin{minipage}[b]{\linewidth}\raggedright
Faixa
\end{minipage} & \begin{minipage}[b]{\linewidth}\raggedright
Tendência
\end{minipage} \\
\midrule\noalign{}
\endhead
\bottomrule\noalign{}
\endlastfoot
Agrícola (com irrigação /\allowbreak  alta pluviosidade) & 1.800 & 1.200 -- 3.000 &
↑ \\
Pastagem estabelecida (pecuária intensiva) & 900 & 600 -- 1.500 & ↑ \\
Terra bruta /\allowbreak  campo nativo & 450 & 200 -- 800 & → \\
\end{longtable}
}

\subsection{Características do
Mercado}

{\def\LTcaptype{none} % do not increment counter
\begin{longtable}[]{@{}ll@{}}
\toprule\noalign{}
Parâmetro & Valor \\
\midrule\noalign{}
\endhead
\bottomrule\noalign{}
\endlastfoot
Valorização anual média (\%) & 5--8\% \\
Lote mínimo rural (ha) & 50 (típico para pecuária viável) \\
Imposto territorial rural (\%) & 1\% sobre valor fiscal \\
Liquidez do mercado & baixa--média \\
\end{longtable}
}
\clearpage
\newpage
\secaoDiagnostico{Boqueron}{\mapaDistrito{1606}}{Boqueron (Boqueron) apresenta perfil operacional consolidado para
comparacao territorial, com leitura conjunta de infraestrutura, risco
hidroclimatico e capacidade de continuidade de servicos essenciais.

\subsubsection{Por que sim}

\begin{itemize}
\tightlist
\item
  Base institucional de referencia disponivel e rastreavel.
\item
  Estrutura comparativa (A-E/\allowbreak GSS) consistente para decisao relativa.
\item
  Plano de mitigacao acionavel ja definido no dossie.
\end{itemize}

\subsubsection{Por que não}

\begin{itemize}
\tightlist
\item
  Serie oficial distrital granular ainda pode apresentar lacunas
  temporais.
\item
  Sensibilidade do score exige monitoramento continuo de risco e
  conectividade.
\item
  Decisao final depende de validacao de campo e janela temporal recente.
\end{itemize}

\subsubsection{Pre-condicoes Minimas de
Decisao}

\begin{itemize}
\tightlist
\item
  Atualizar indicadores criticos em ciclo trimestral.
\item
  Confirmar trafegabilidade e contingencia hidrica/\allowbreak energetica local.
\item
  Validar checklist de resiliencia domiciliar/\allowbreak logistica antes de decisao
  final.
\end{itemize}

\subsubsection{Recomendação
Técnica}

Uso recomendado como candidato em analise multicriterio, condicionado ao
cumprimento das pre-condicoes acima. Referência atual de score: GSS 7.7;
risco predominante: moderado.

\subsubsection{}

\begin{itemize}
\tightlist
\item
  \begin{enumerate}
  \def\labelenumi{\arabic{enumi})}
  \tightlist
  \item
    Delimitacao territorial validada por georreferencia institucional.
  \end{enumerate}
\item
  \begin{enumerate}
  \def\labelenumi{\arabic{enumi})}
  \setcounter{enumi}{1}
  \tightlist
  \item
    População municipal/\allowbreak distrital referenciada por base censitaria
    nacional.
  \end{enumerate}
\item
  \begin{enumerate}
  \def\labelenumi{\arabic{enumi})}
  \setcounter{enumi}{2}
  \tightlist
  \item
    Comparabilidade intra-departamental apoiada em indicadores
    distritais oficiais.
  \end{enumerate}
\item
  \begin{enumerate}
  \def\labelenumi{\arabic{enumi})}
  \setcounter{enumi}{3}
  \tightlist
  \item
    Estado macro de conectividade viaria ancorado em informacoes de
    rodovias.
  \end{enumerate}
\item
  \begin{enumerate}
  \def\labelenumi{\arabic{enumi})}
  \setcounter{enumi}{4}
  \tightlist
  \item
    Exposicao hidrometeorologica monitorada por avisos oficiais.
  \end{enumerate}
\item
  \begin{enumerate}
  \def\labelenumi{\arabic{enumi})}
  \setcounter{enumi}{5}
  \tightlist
  \item
    Historico de contexto meteorologico apoiado em publicacoes tecnicas.
  \end{enumerate}
\item
  \begin{enumerate}
  \def\labelenumi{\arabic{enumi})}
  \setcounter{enumi}{6}
  \tightlist
  \item
    Capacidade institucional de resposta a eventos apoiada em acoes da
    SEN\footnote{Secretaría de Emergencia Nacional (Secretaria de Emergência Nacional), órgão governamental de resposta a desastres.}.
  \end{enumerate}
\item
  \begin{enumerate}
  \def\labelenumi{\arabic{enumi})}
  \setcounter{enumi}{7}
  \tightlist
  \item
    Infraestrutura energetica analisada via operador nacional.
  \end{enumerate}
\item
  \begin{enumerate}
  \def\labelenumi{\arabic{enumi})}
  \setcounter{enumi}{8}
  \tightlist
  \item
    Referências de obras e logistica nacional verificadas no MOPC.
  \end{enumerate}
\item
  \begin{enumerate}
  \def\labelenumi{\arabic{enumi})}
  \setcounter{enumi}{9}
  \tightlist
  \item
    Contexto sociopolitico institucional referenciado por autoridade
    eleitoral.
  \end{enumerate}
\item
  \begin{enumerate}
  \def\labelenumi{\arabic{enumi})}
  \setcounter{enumi}{10}
  \tightlist
  \item
    Camada cartografica de apoio eleitoral/\allowbreak territorial consultada.
  \end{enumerate}
\item
  \begin{enumerate}
  \def\labelenumi{\arabic{enumi})}
  \setcounter{enumi}{11}
  \tightlist
  \item
    Dados publicos complementares para verificacao cruzada.
  \end{enumerate}
\end{itemize}}{dist:16_Boqueron:Boqueron}
\begin{table}[ht!]
\small \begin{tabular}{p{3.0cm}p{0.8cm}p{6.0cm}} \toprule
\textbf{Critério} & \textbf{Nota} & \textbf{Justificativa} \\ \midrule
A - Ameaças Estratégicas & 7.8 & - \\
B - Risco Social & 8.4 & - \\
C - Riscos Naturais & 6.4 & - \\
D - Autossuficiência & 7.5 & - \\
E - Institucional & 8.2 & - \\
\midrule \textbf{GSS FINAL} & \textbf{7.7} & \textbf{Seguro (bom potencial com preparacao).} \\ \bottomrule \end{tabular}
\end{table}
\subsubsection{Vulnerabilidades e
Mitigação}\label{vuln-Boqueron-vulnerabilidades-e-mitigauxe7uxe3o}

\begin{itemize}
\tightlist
\item
  \textbf{Isolamento Geográfico:} Dependência de transporte fluvial e
  estradas precárias (Mitigação: manter embarcação própria e estoque de
  suprimentos para 30 dias).
\item
  \textbf{Risco de Inundação:} Vulnerabilidade extrema a cheias do Rio
  Paraguai (Mitigação: construção em cotas elevadas, evitar áreas
  ribeirinhas diretas).
\item
  \textbf{Instabilidade Energética:} Final de linha de transmissão
  (Mitigação: sistemas solares de backup com baterias).
\item
  \textbf{Segurança Regional:} Monitoramento de atividades de grupos
  armados no departamento.
\end{itemize}

\subsubsection{Dossiê de Campo}\label{dossie-Boqueron-dossiuxea-de-campo}
\begin{description}[style=nextline,leftmargin=0.5cm,font=\bfseries]
\item[Coordenadas]  22°59'S 57°43'W.
\item[Topografia]  Margem esquerda do Rio Paraguai. Relevo predominantemente plano com formações calcárias. Geologicamente estável, risco sísmico muito baixo.
\item[Alvos Estratégicos]  Indústrias de calcário (Calerías), Porto de Itacuá (escoamento de cimento e minerais), proximidade com a fronteira brasileira via rio.
\item[Fallout]  Ventos predominantemente do Norte/\allowbreak Nordeste (quentes/\allowbreak úmidos) e Sul (frentes frias).
\item[População]  2.377 (Censo 2022 Final).
\item[Densidade]  \textasciitilde 1,39 hab/\allowbreak km² (Extremamente baixa, excelente para isolamento estratégico).
\item[Vias de Acesso]  Acesso terrestre precário (Ruta PY22 e ramais de terra), forte dependência de transporte fluvial. Pista de pouso asfaltada em Vallemí (\textasciitilde 10 km).
\item[Serviços]  Infraestrutura básica local. Dependência de Concepción (190 km) ou Pedro Juan Caballero para saúde de alta complexidade.
\item[Custo de Vida]  \textasciitilde 30-40\% mais baixo que centros urbanos brasileiros; energia e água baratas.
\item[Preço da Terra]  US\$ 1.500-5.000/\allowbreak ha (rural/\allowbreak pastagem); lotes urbanos em San Lázaro por US\$ 15.000-30.000.
\item[Hidrologia]  Risco Crítico de Inundação. Proximidade direta com o Rio Paraguai; histórico de cheias severas (2019, 2023) que podem isolar a localidade por terra.
\end{description}
\textbf{Fonte climática:} NASA POWER Climatology API (período 2001-2020)

\textbf{Fonte luminosa:} estimativa world atlas

\paragraph{Irradiação Solar
(kWh/\allowbreak m²/\allowbreak dia)}\label{irradiauxe7uxe3o-solar-kwhmuxb2dia}

{\def\LTcaptype{none} % do not increment counter
\begin{longtable}[]{@{}
  >{\raggedright\arraybackslash}p{(\linewidth - 24\tabcolsep) * \real{0.0746}}
  >{\raggedright\arraybackslash}p{(\linewidth - 24\tabcolsep) * \real{0.0746}}
  >{\raggedright\arraybackslash}p{(\linewidth - 24\tabcolsep) * \real{0.0746}}
  >{\raggedright\arraybackslash}p{(\linewidth - 24\tabcolsep) * \real{0.0746}}
  >{\raggedright\arraybackslash}p{(\linewidth - 24\tabcolsep) * \real{0.0746}}
  >{\raggedright\arraybackslash}p{(\linewidth - 24\tabcolsep) * \real{0.0746}}
  >{\raggedright\arraybackslash}p{(\linewidth - 24\tabcolsep) * \real{0.0746}}
  >{\raggedright\arraybackslash}p{(\linewidth - 24\tabcolsep) * \real{0.0746}}
  >{\raggedright\arraybackslash}p{(\linewidth - 24\tabcolsep) * \real{0.0746}}
  >{\raggedright\arraybackslash}p{(\linewidth - 24\tabcolsep) * \real{0.0746}}
  >{\raggedright\arraybackslash}p{(\linewidth - 24\tabcolsep) * \real{0.0746}}
  >{\raggedright\arraybackslash}p{(\linewidth - 24\tabcolsep) * \real{0.0746}}
  >{\raggedright\arraybackslash}p{(\linewidth - 24\tabcolsep) * \real{0.1045}}@{}}
\toprule\noalign{}
\begin{minipage}[b]{\linewidth}\raggedright
Jan
\end{minipage} & \begin{minipage}[b]{\linewidth}\raggedright
Fev
\end{minipage} & \begin{minipage}[b]{\linewidth}\raggedright
Mar
\end{minipage} & \begin{minipage}[b]{\linewidth}\raggedright
Abr
\end{minipage} & \begin{minipage}[b]{\linewidth}\raggedright
Mai
\end{minipage} & \begin{minipage}[b]{\linewidth}\raggedright
Jun
\end{minipage} & \begin{minipage}[b]{\linewidth}\raggedright
Jul
\end{minipage} & \begin{minipage}[b]{\linewidth}\raggedright
Ago
\end{minipage} & \begin{minipage}[b]{\linewidth}\raggedright
Set
\end{minipage} & \begin{minipage}[b]{\linewidth}\raggedright
Out
\end{minipage} & \begin{minipage}[b]{\linewidth}\raggedright
Nov
\end{minipage} & \begin{minipage}[b]{\linewidth}\raggedright
Dez
\end{minipage} & \begin{minipage}[b]{\linewidth}\raggedright
Média
\end{minipage} \\
\midrule\noalign{}
\endhead
\bottomrule\noalign{}
\endlastfoot
6.55 & 5.94 & 5.63 & 4.73 & 3.58 & 3.17 & 3.41 & 4.19 & 4.70 & 5.44 &
6.28 & 6.49 & \textbf{5.01} \\
\end{longtable}
}

\textbf{Inclinação solar recomendada:} 22° N (anual) · 32° N (inverno
jun-ago) · 12° N (verão nov-jan)

\paragraph{Precipitação (mm/\allowbreak mês)}\label{precipitauxe7uxe3o-mmmuxeas}

{\def\LTcaptype{none} % do not increment counter
\begin{longtable}[]{@{}
  >{\raggedright\arraybackslash}p{(\linewidth - 24\tabcolsep) * \real{0.0704}}
  >{\raggedright\arraybackslash}p{(\linewidth - 24\tabcolsep) * \real{0.0704}}
  >{\raggedright\arraybackslash}p{(\linewidth - 24\tabcolsep) * \real{0.0704}}
  >{\raggedright\arraybackslash}p{(\linewidth - 24\tabcolsep) * \real{0.0704}}
  >{\raggedright\arraybackslash}p{(\linewidth - 24\tabcolsep) * \real{0.0704}}
  >{\raggedright\arraybackslash}p{(\linewidth - 24\tabcolsep) * \real{0.0704}}
  >{\raggedright\arraybackslash}p{(\linewidth - 24\tabcolsep) * \real{0.0704}}
  >{\raggedright\arraybackslash}p{(\linewidth - 24\tabcolsep) * \real{0.0704}}
  >{\raggedright\arraybackslash}p{(\linewidth - 24\tabcolsep) * \real{0.0704}}
  >{\raggedright\arraybackslash}p{(\linewidth - 24\tabcolsep) * \real{0.0704}}
  >{\raggedright\arraybackslash}p{(\linewidth - 24\tabcolsep) * \real{0.0704}}
  >{\raggedright\arraybackslash}p{(\linewidth - 24\tabcolsep) * \real{0.0704}}
  >{\raggedright\arraybackslash}p{(\linewidth - 24\tabcolsep) * \real{0.1549}}@{}}
\toprule\noalign{}
\begin{minipage}[b]{\linewidth}\raggedright
Jan
\end{minipage} & \begin{minipage}[b]{\linewidth}\raggedright
Fev
\end{minipage} & \begin{minipage}[b]{\linewidth}\raggedright
Mar
\end{minipage} & \begin{minipage}[b]{\linewidth}\raggedright
Abr
\end{minipage} & \begin{minipage}[b]{\linewidth}\raggedright
Mai
\end{minipage} & \begin{minipage}[b]{\linewidth}\raggedright
Jun
\end{minipage} & \begin{minipage}[b]{\linewidth}\raggedright
Jul
\end{minipage} & \begin{minipage}[b]{\linewidth}\raggedright
Ago
\end{minipage} & \begin{minipage}[b]{\linewidth}\raggedright
Set
\end{minipage} & \begin{minipage}[b]{\linewidth}\raggedright
Out
\end{minipage} & \begin{minipage}[b]{\linewidth}\raggedright
Nov
\end{minipage} & \begin{minipage}[b]{\linewidth}\raggedright
Dez
\end{minipage} & \begin{minipage}[b]{\linewidth}\raggedright
Total/\allowbreak ano
\end{minipage} \\
\midrule\noalign{}
\endhead
\bottomrule\noalign{}
\endlastfoot
152 & 149 & 120 & 118 & 106 & 57 & 38 & 23 & 55 & 130 & 202 & 167 &
\textbf{1321 mm} \\
\end{longtable}
}

\paragraph{Poluição Luminosa}\label{poluiuxe7uxe3o-luminosa}

{\def\LTcaptype{none} % do not increment counter
\begin{longtable}[]{@{}ll@{}}
\toprule\noalign{}
Parâmetro & Valor \\
\midrule\noalign{}
\endhead
\bottomrule\noalign{}
\endlastfoot
Escala Bortle & 3 --- Céu rural \\
Radiância artificial & 1.5 nW/\allowbreak cm²/\allowbreak sr \\
\end{longtable}
}
\paragraph{Solo (SoilGrids 2.0, média ponderada 0--30
cm)}

{\def\LTcaptype{none} % do not increment counter
\begin{longtable}[]{@{}ll@{}}
\toprule\noalign{}
Parâmetro & Valor \\
\midrule\noalign{}
\endhead
\bottomrule\noalign{}
\endlastfoot
pH (H$_2$O) & 5.95 \\
Carbono orgânico (SOC) & 17.27 g/\allowbreak kg \\
Argila & 21.87 \% \\
Areia & 39.29 \% \\
Silte (calc.) & 38.8 \% \\
Densidade aparente & 1.33 g/\allowbreak cm³ \\
\textbf{Aptidão agrícola} & \textbf{Alta} \\
\end{longtable}
}

\subsubsection{Indicadores Sociais}

\textbf{Fontes:} PNUD 2020, INE\footnote{Instituto Nacional de Estadística (Instituto Nacional de Estatística), órgão oficial de dados demográficos e censitários.} Censo 2022, INE\footnote{Instituto Nacional de Estadística (Instituto Nacional de Estatística), órgão oficial de dados demográficos e censitários.} EPHC 2023, Ministerio
Público 2024\\
\textbf{Nota:} Valores marcados como \emph{(dept.)} referem-se ao
departamento de Concepción; valores \emph{(dist.)} são específicos deste
distrito. Dados marcados \emph{(est.)} são estimativas.

{\def\LTcaptype{none} % do not increment counter
\begin{longtable}[]{@{}
  >{\raggedright\arraybackslash}p{(\linewidth - 4\tabcolsep) * \real{0.4231}}
  >{\raggedright\arraybackslash}p{(\linewidth - 4\tabcolsep) * \real{0.2692}}
  >{\raggedright\arraybackslash}p{(\linewidth - 4\tabcolsep) * \real{0.3077}}@{}}
\toprule\noalign{}
\begin{minipage}[b]{\linewidth}\raggedright
Indicador
\end{minipage} & \begin{minipage}[b]{\linewidth}\raggedright
Valor
\end{minipage} & \begin{minipage}[b]{\linewidth}\raggedright
Âmbito
\end{minipage} \\
\midrule\noalign{}
\endhead
\bottomrule\noalign{}
\endlastfoot
IDH (2020) & 0.630 & dist. (est.) \\
Ranking IDH nacional & 9/\allowbreak 18 & dept. \\
Esperança de vida & 72,4 anos & dept. \\
Escolaridade média & 6.9 anos & dist. \\
RNB per capita (USD PPA) & 8.400 & dept. \\
Pobreza monetária (\%) & 31,1\% & dept. \\
Pobreza extrema (\%) & 7,8\% & dept. \\
Índice de Gini & 0,460 & dept. \\
Acesso a água potável (\%) & 74,2\% & dept. \\
Acesso a saneamento (\%) & 71,8\% & dept. \\
Taxa de homicídios (est., /\allowbreak 100k hab) & \textasciitilde12--15
(estimativa, acima da média nacional) & dept. \\
Índice de segurança & baixo & dept. \\
População (Censo 2022) & 2.377 hab. & dist. \\
Idade mediana & 25.0 anos & dist. \\
Taxa de fecundidade & 3 & dist. \\
\end{longtable}
}
\subsubsection{Combustível}

\textbf{Referência:} PETROPAR /\allowbreak  postos locais (2024) \textbf{Tipo de
localidade:} Interior

{\def\LTcaptype{none} % do not increment counter
\begin{longtable}[]{@{}lll@{}}
\toprule\noalign{}
Combustível & USD/\allowbreak litro & Gs/\allowbreak litro (aprox.) \\
\midrule\noalign{}
\endhead
\bottomrule\noalign{}
\endlastfoot
Gasolina 93 oct & 1.01 & 7,474 \\
Gasolina 97 oct (premium) & 1.11 & 8,214 \\
Diesel & 0.94 & 6,956 \\
\end{longtable}
}

\begin{quote}
Preços podem variar ±5\% conforme posto e sazonalidade. Chaco e interior
remoto apresentam maior variação.
\end{quote}

\subsubsection{Cobertura Celular}

\textbf{Fonte:} CONATEL PY /\allowbreak  operadoras (2024)

{\def\LTcaptype{none} % do not increment counter
\begin{longtable}[]{@{}ll@{}}
\toprule\noalign{}
Parâmetro & Valor \\
\midrule\noalign{}
\endhead
\bottomrule\noalign{}
\endlastfoot
Cobertura 4G população (dept.) & 78\% \\
Cobertura 4G área rural & 55\% \\
Melhor operadora & Tigo \\
Qualidade rural & moderada \\
\end{longtable}
}

\begin{quote}
Para áreas rurais fora do núcleo urbano, recomenda-se chip Tigo como
principal e Personal como backup.
\end{quote}

\subsubsection{Internet}

\textbf{Fonte:} CONATEL /\allowbreak  Speedtest Ookla (2024)

{\def\LTcaptype{none} % do not increment counter
\begin{longtable}[]{@{}ll@{}}
\toprule\noalign{}
Parâmetro & Valor \\
\midrule\noalign{}
\endhead
\bottomrule\noalign{}
\endlastfoot
Velocidade média download & 38 Mbps \\
Domicílios com internet (dept.) & 48\% \\
Tecnologia predominante & rádio \\
Opção rural & Starlink disponível (\textasciitilde USD 44/\allowbreak mês) \\
\end{longtable}
}

\subsubsection{Mercado Imobiliário e Terra
Rural}

\textbf{Fonte:} INDERT /\allowbreak  Clasificados.com.py (2024)

{\def\LTcaptype{none} % do not increment counter
\begin{longtable}[]{@{}ll@{}}
\toprule\noalign{}
Tipo & Referência \\
\midrule\noalign{}
\endhead
\bottomrule\noalign{}
\endlastfoot
Terra agrícola alta prod. (USD/\allowbreak ha) & 3,500 \\
Imóvel urbano (USD/\allowbreak m²) & 650 \\
Aluguel 2 quartos (USD/\allowbreak mês) & 260 \\
\end{longtable}
}

\begin{quote}
Valores de referência departamental. Localidades menores podem ter
preços 20--40\% abaixo da capital departamental.
\end{quote}

\subsubsection{Saúde}

\textbf{Fonte:} MSPBS /\allowbreak  IPS Paraguay (2024-2026), consolidação
departamental e proxy local

{\def\LTcaptype{none} % do not increment counter
\begin{longtable}[]{@{}ll@{}}
\toprule\noalign{}
Serviço & Disponibilidade \\
\midrule\noalign{}
\endhead
\bottomrule\noalign{}
\endlastfoot
USF /\allowbreak  Posto de Saúde & sim \\
Hospital Regional & não \\
IPS (seguro social) & sim \\
Farmácia & sim \\
Distância ao hospital de referência & \textasciitilde138 km
(Concepción) \\
\end{longtable}
}

\textbf{Principais estabelecimentos:} USF local; referencia hospitalar
em Concepción

\textbf{Observação para imigrantes:} Atendimento primario local ou em
raio curto; casos de maior complexidade seguem para Concepción.
Cobertura privada continua recomendada para especialidades e urgências
de maior complexidade.
\newpage
\secaoDiagnostico{Filadelfia}{\mapaDistrito{1604}}{Loreto (Concepción) apresenta perfil operacional consolidado para
comparacao territorial, com leitura conjunta de infraestrutura, risco
hidroclimatico e capacidade de continuidade de servicos essenciais.

\subsubsection{Por que sim}

\begin{itemize}
\tightlist
\item
  Base institucional de referencia disponivel e rastreavel.
\item
  Estrutura comparativa (A-E/\allowbreak GSS) consistente para decisao relativa.
\item
  Plano de mitigacao acionavel ja definido no dossie.
\end{itemize}

\subsubsection{Por que não}

\begin{itemize}
\tightlist
\item
  Serie oficial distrital granular ainda pode apresentar lacunas
  temporais.
\item
  Sensibilidade do score exige monitoramento continuo de risco e
  conectividade.
\item
  Decisao final depende de validacao de campo e janela temporal recente.
\end{itemize}

\subsubsection{Pre-condicoes Minimas de
Decisao}

\begin{itemize}
\tightlist
\item
  Atualizar indicadores criticos em ciclo trimestral.
\item
  Confirmar trafegabilidade e contingencia hidrica/\allowbreak energetica local.
\item
  Validar checklist de resiliencia domiciliar/\allowbreak logistica antes de decisao
  final.
\end{itemize}

\subsubsection{Recomendação
Técnica}

Uso recomendado como candidato em analise multicriterio, condicionado ao
cumprimento das pre-condicoes acima. Referência atual de score: GSS 7.9;
risco predominante: baixo.

\subsubsection{}

\begin{itemize}
\tightlist
\item
  \begin{enumerate}
  \def\labelenumi{\arabic{enumi})}
  \tightlist
  \item
    Delimitacao territorial validada por georreferencia institucional.
  \end{enumerate}
\item
  \begin{enumerate}
  \def\labelenumi{\arabic{enumi})}
  \setcounter{enumi}{1}
  \tightlist
  \item
    População municipal/\allowbreak distrital referenciada por base censitaria
    nacional.
  \end{enumerate}
\item
  \begin{enumerate}
  \def\labelenumi{\arabic{enumi})}
  \setcounter{enumi}{2}
  \tightlist
  \item
    Comparabilidade intra-departamental apoiada em indicadores
    distritais oficiais.
  \end{enumerate}
\item
  \begin{enumerate}
  \def\labelenumi{\arabic{enumi})}
  \setcounter{enumi}{3}
  \tightlist
  \item
    Estado macro de conectividade viaria ancorado em informacoes de
    rodovias.
  \end{enumerate}
\item
  \begin{enumerate}
  \def\labelenumi{\arabic{enumi})}
  \setcounter{enumi}{4}
  \tightlist
  \item
    Exposicao hidrometeorologica monitorada por avisos oficiais.
  \end{enumerate}
\item
  \begin{enumerate}
  \def\labelenumi{\arabic{enumi})}
  \setcounter{enumi}{5}
  \tightlist
  \item
    Historico de contexto meteorologico apoiado em publicacoes tecnicas.
  \end{enumerate}
\item
  \begin{enumerate}
  \def\labelenumi{\arabic{enumi})}
  \setcounter{enumi}{6}
  \tightlist
  \item
    Capacidade institucional de resposta a eventos apoiada em acoes da
    SEN\footnote{Secretaría de Emergencia Nacional (Secretaria de Emergência Nacional), órgão governamental de resposta a desastres.}.
  \end{enumerate}
\item
  \begin{enumerate}
  \def\labelenumi{\arabic{enumi})}
  \setcounter{enumi}{7}
  \tightlist
  \item
    Infraestrutura energetica analisada via operador nacional.
  \end{enumerate}
\item
  \begin{enumerate}
  \def\labelenumi{\arabic{enumi})}
  \setcounter{enumi}{8}
  \tightlist
  \item
    Referências de obras e logistica nacional verificadas no MOPC.
  \end{enumerate}
\item
  \begin{enumerate}
  \def\labelenumi{\arabic{enumi})}
  \setcounter{enumi}{9}
  \tightlist
  \item
    Contexto sociopolitico institucional referenciado por autoridade
    eleitoral.
  \end{enumerate}
\item
  \begin{enumerate}
  \def\labelenumi{\arabic{enumi})}
  \setcounter{enumi}{10}
  \tightlist
  \item
    Camada cartografica de apoio eleitoral/\allowbreak territorial consultada.
  \end{enumerate}
\item
  \begin{enumerate}
  \def\labelenumi{\arabic{enumi})}
  \setcounter{enumi}{11}
  \tightlist
  \item
    Dados publicos complementares para verificacao cruzada.
  \end{enumerate}
\end{itemize}}{dist:16_Boqueron:Filadelfia}
\begin{table}[ht!]
\small \begin{tabular}{p{3.0cm}p{0.8cm}p{6.0cm}} \toprule
\textbf{Critério} & \textbf{Nota} & \textbf{Justificativa} \\ \midrule
A - Ameaças Estratégicas & 7.6 & - \\
B - Risco Social & 8.3 & - \\
C - Riscos Naturais & 6.5 & - \\
D - Autossuficiência & 8.0 & - \\
E - Institucional & 8.2 & - \\
\midrule \textbf{GSS FINAL} & \textbf{7.9} & \textbf{Seguro (bom potencial com preparacao).} \\ \bottomrule \end{tabular}
\end{table}
\subsubsection{Vulnerabilidades e
Mitigação}\label{vuln-Filadelfia-vulnerabilidades-e-mitigauxe7uxe3o}

\begin{itemize}
\tightlist
\item
  Exposicao hidrometeorologica controlada (\textless45/\allowbreak 100).
\item
  Conectividade viaria favoravel (\textgreater=70/\allowbreak 100).
\item
  Estoque de contingencia para 14 dias (agua/\allowbreak alimentos/\allowbreak combustivel),
  priorizado quando risco hidro=44/\allowbreak 100.
\item
  Duplicar rotas/\allowbreak logistica quando conectividade viaria=85/\allowbreak 100 for
  \textless70; manter rota secundaria mapeada e testada trimestralmente.
\item
  Plano de energia resiliente com autonomia minima de 72h
  (geracao/\allowbreak backup), revisao semestral.
\end{itemize}

Sintese aplicada: - Dados textuais consolidados com base nas fontes
oficiais acima; proximos ciclos podem aprofundar indicadores numericos
especificos por distrito.
\subsubsection{Dossiê de Campo}\label{dossie-Filadelfia-dossiuxea-de-campo}
\begin{description}[style=nextline,leftmargin=0.5cm,font=\bfseries]
\item[Coordenadas]  23°16'S 57°19'W (geocodificado via Nominatim).
\end{description}
\textbf{Fonte climática:} NASA POWER Climatology API (período 2001-2020)

\textbf{Fonte luminosa:} estimativa world atlas

\paragraph{Irradiação Solar
(kWh/\allowbreak m²/\allowbreak dia)}\label{irradiauxe7uxe3o-solar-kwhmuxb2dia}

{\def\LTcaptype{none} % do not increment counter
\begin{longtable}[]{@{}
  >{\raggedright\arraybackslash}p{(\linewidth - 24\tabcolsep) * \real{0.0746}}
  >{\raggedright\arraybackslash}p{(\linewidth - 24\tabcolsep) * \real{0.0746}}
  >{\raggedright\arraybackslash}p{(\linewidth - 24\tabcolsep) * \real{0.0746}}
  >{\raggedright\arraybackslash}p{(\linewidth - 24\tabcolsep) * \real{0.0746}}
  >{\raggedright\arraybackslash}p{(\linewidth - 24\tabcolsep) * \real{0.0746}}
  >{\raggedright\arraybackslash}p{(\linewidth - 24\tabcolsep) * \real{0.0746}}
  >{\raggedright\arraybackslash}p{(\linewidth - 24\tabcolsep) * \real{0.0746}}
  >{\raggedright\arraybackslash}p{(\linewidth - 24\tabcolsep) * \real{0.0746}}
  >{\raggedright\arraybackslash}p{(\linewidth - 24\tabcolsep) * \real{0.0746}}
  >{\raggedright\arraybackslash}p{(\linewidth - 24\tabcolsep) * \real{0.0746}}
  >{\raggedright\arraybackslash}p{(\linewidth - 24\tabcolsep) * \real{0.0746}}
  >{\raggedright\arraybackslash}p{(\linewidth - 24\tabcolsep) * \real{0.0746}}
  >{\raggedright\arraybackslash}p{(\linewidth - 24\tabcolsep) * \real{0.1045}}@{}}
\toprule\noalign{}
\begin{minipage}[b]{\linewidth}\raggedright
Jan
\end{minipage} & \begin{minipage}[b]{\linewidth}\raggedright
Fev
\end{minipage} & \begin{minipage}[b]{\linewidth}\raggedright
Mar
\end{minipage} & \begin{minipage}[b]{\linewidth}\raggedright
Abr
\end{minipage} & \begin{minipage}[b]{\linewidth}\raggedright
Mai
\end{minipage} & \begin{minipage}[b]{\linewidth}\raggedright
Jun
\end{minipage} & \begin{minipage}[b]{\linewidth}\raggedright
Jul
\end{minipage} & \begin{minipage}[b]{\linewidth}\raggedright
Ago
\end{minipage} & \begin{minipage}[b]{\linewidth}\raggedright
Set
\end{minipage} & \begin{minipage}[b]{\linewidth}\raggedright
Out
\end{minipage} & \begin{minipage}[b]{\linewidth}\raggedright
Nov
\end{minipage} & \begin{minipage}[b]{\linewidth}\raggedright
Dez
\end{minipage} & \begin{minipage}[b]{\linewidth}\raggedright
Média
\end{minipage} \\
\midrule\noalign{}
\endhead
\bottomrule\noalign{}
\endlastfoot
6.68 & 6.07 & 5.58 & 4.63 & 3.47 & 3.00 & 3.30 & 4.09 & 4.68 & 5.49 &
6.34 & 6.61 & \textbf{5.0} \\
\end{longtable}
}

\textbf{Inclinação solar recomendada:} 23° N (anual) · 33° N (inverno
jun-ago) · 13° N (verão nov-jan)

\paragraph{Precipitação (mm/\allowbreak mês)}\label{precipitauxe7uxe3o-mmmuxeas}

{\def\LTcaptype{none} % do not increment counter
\begin{longtable}[]{@{}
  >{\raggedright\arraybackslash}p{(\linewidth - 24\tabcolsep) * \real{0.0704}}
  >{\raggedright\arraybackslash}p{(\linewidth - 24\tabcolsep) * \real{0.0704}}
  >{\raggedright\arraybackslash}p{(\linewidth - 24\tabcolsep) * \real{0.0704}}
  >{\raggedright\arraybackslash}p{(\linewidth - 24\tabcolsep) * \real{0.0704}}
  >{\raggedright\arraybackslash}p{(\linewidth - 24\tabcolsep) * \real{0.0704}}
  >{\raggedright\arraybackslash}p{(\linewidth - 24\tabcolsep) * \real{0.0704}}
  >{\raggedright\arraybackslash}p{(\linewidth - 24\tabcolsep) * \real{0.0704}}
  >{\raggedright\arraybackslash}p{(\linewidth - 24\tabcolsep) * \real{0.0704}}
  >{\raggedright\arraybackslash}p{(\linewidth - 24\tabcolsep) * \real{0.0704}}
  >{\raggedright\arraybackslash}p{(\linewidth - 24\tabcolsep) * \real{0.0704}}
  >{\raggedright\arraybackslash}p{(\linewidth - 24\tabcolsep) * \real{0.0704}}
  >{\raggedright\arraybackslash}p{(\linewidth - 24\tabcolsep) * \real{0.0704}}
  >{\raggedright\arraybackslash}p{(\linewidth - 24\tabcolsep) * \real{0.1549}}@{}}
\toprule\noalign{}
\begin{minipage}[b]{\linewidth}\raggedright
Jan
\end{minipage} & \begin{minipage}[b]{\linewidth}\raggedright
Fev
\end{minipage} & \begin{minipage}[b]{\linewidth}\raggedright
Mar
\end{minipage} & \begin{minipage}[b]{\linewidth}\raggedright
Abr
\end{minipage} & \begin{minipage}[b]{\linewidth}\raggedright
Mai
\end{minipage} & \begin{minipage}[b]{\linewidth}\raggedright
Jun
\end{minipage} & \begin{minipage}[b]{\linewidth}\raggedright
Jul
\end{minipage} & \begin{minipage}[b]{\linewidth}\raggedright
Ago
\end{minipage} & \begin{minipage}[b]{\linewidth}\raggedright
Set
\end{minipage} & \begin{minipage}[b]{\linewidth}\raggedright
Out
\end{minipage} & \begin{minipage}[b]{\linewidth}\raggedright
Nov
\end{minipage} & \begin{minipage}[b]{\linewidth}\raggedright
Dez
\end{minipage} & \begin{minipage}[b]{\linewidth}\raggedright
Total/\allowbreak ano
\end{minipage} \\
\midrule\noalign{}
\endhead
\bottomrule\noalign{}
\endlastfoot
157 & 151 & 120 & 139 & 116 & 71 & 48 & 31 & 73 & 151 & 204 & 167 &
\textbf{1432 mm} \\
\end{longtable}
}

\paragraph{Poluição Luminosa}\label{poluiuxe7uxe3o-luminosa}

{\def\LTcaptype{none} % do not increment counter
\begin{longtable}[]{@{}ll@{}}
\toprule\noalign{}
Parâmetro & Valor \\
\midrule\noalign{}
\endhead
\bottomrule\noalign{}
\endlastfoot
Escala Bortle & 3 --- Céu rural \\
Radiância artificial & 1.5 nW/\allowbreak cm²/\allowbreak sr \\
\end{longtable}
}
\paragraph{Solo (SoilGrids 2.0, média ponderada 0--30
cm)}

{\def\LTcaptype{none} % do not increment counter
\begin{longtable}[]{@{}ll@{}}
\toprule\noalign{}
Parâmetro & Valor \\
\midrule\noalign{}
\endhead
\bottomrule\noalign{}
\endlastfoot
pH (H$_2$O) & 5.75 \\
Carbono orgânico (SOC) & 14.28 g/\allowbreak kg \\
Argila & 18.2 \% \\
Areia & 61.83 \% \\
Silte (calc.) & 20.0 \% \\
Densidade aparente & 1.35 g/\allowbreak cm³ \\
\textbf{Aptidão agrícola} & \textbf{Média} \\
\end{longtable}
}

\subsubsection{Indicadores Sociais}

\textbf{Fontes:} PNUD 2020, INE\footnote{Instituto Nacional de Estadística (Instituto Nacional de Estatística), órgão oficial de dados demográficos e censitários.} Censo 2022, INE\footnote{Instituto Nacional de Estadística (Instituto Nacional de Estatística), órgão oficial de dados demográficos e censitários.} EPHC 2023, Ministerio
Público 2024\\
\textbf{Nota:} Valores marcados como \emph{(dept.)} referem-se ao
departamento de Concepción; valores \emph{(dist.)} são específicos deste
distrito. Dados marcados \emph{(est.)} são estimativas.

{\def\LTcaptype{none} % do not increment counter
\begin{longtable}[]{@{}
  >{\raggedright\arraybackslash}p{(\linewidth - 4\tabcolsep) * \real{0.4231}}
  >{\raggedright\arraybackslash}p{(\linewidth - 4\tabcolsep) * \real{0.2692}}
  >{\raggedright\arraybackslash}p{(\linewidth - 4\tabcolsep) * \real{0.3077}}@{}}
\toprule\noalign{}
\begin{minipage}[b]{\linewidth}\raggedright
Indicador
\end{minipage} & \begin{minipage}[b]{\linewidth}\raggedright
Valor
\end{minipage} & \begin{minipage}[b]{\linewidth}\raggedright
Âmbito
\end{minipage} \\
\midrule\noalign{}
\endhead
\bottomrule\noalign{}
\endlastfoot
IDH (2020) & 0.650 & dist. (est.) \\
Ranking IDH nacional & 9/\allowbreak 18 & dept. \\
Esperança de vida & 72,4 anos & dept. \\
Escolaridade média & 8.7 anos & dist. \\
RNB per capita (USD PPA) & 8.400 & dept. \\
Pobreza monetária (\%) & 31,1\% & dept. \\
Pobreza extrema (\%) & 7,8\% & dept. \\
Índice de Gini & 0,460 & dept. \\
Acesso a água potável (\%) & 74,2\% & dept. \\
Acesso a saneamento (\%) & 71,8\% & dept. \\
Taxa de homicídios (est., /\allowbreak 100k hab) & \textasciitilde12--15
(estimativa, acima da média nacional) & dept. \\
Índice de segurança & baixo & dept. \\
População (Censo 2022) & 1.358 hab. & dist. \\
Idade mediana & 29.0 anos & dist. \\
Taxa de fecundidade & 2 & dist. \\
\end{longtable}
}
\subsubsection{Combustível}

\textbf{Referência:} PETROPAR /\allowbreak  postos locais (2024) \textbf{Tipo de
localidade:} Interior

{\def\LTcaptype{none} % do not increment counter
\begin{longtable}[]{@{}lll@{}}
\toprule\noalign{}
Combustível & USD/\allowbreak litro & Gs/\allowbreak litro (aprox.) \\
\midrule\noalign{}
\endhead
\bottomrule\noalign{}
\endlastfoot
Gasolina 93 oct & 1.01 & 7,474 \\
Gasolina 97 oct (premium) & 1.11 & 8,214 \\
Diesel & 0.94 & 6,956 \\
\end{longtable}
}

\begin{quote}
Preços podem variar ±5\% conforme posto e sazonalidade. Chaco e interior
remoto apresentam maior variação.
\end{quote}

\subsubsection{Cobertura Celular}

\textbf{Fonte:} CONATEL PY /\allowbreak  operadoras (2024)

{\def\LTcaptype{none} % do not increment counter
\begin{longtable}[]{@{}ll@{}}
\toprule\noalign{}
Parâmetro & Valor \\
\midrule\noalign{}
\endhead
\bottomrule\noalign{}
\endlastfoot
Cobertura 4G população (dept.) & 78\% \\
Cobertura 4G área rural & 55\% \\
Melhor operadora & Tigo \\
Qualidade rural & moderada \\
\end{longtable}
}

\begin{quote}
Para áreas rurais fora do núcleo urbano, recomenda-se chip Tigo como
principal e Personal como backup.
\end{quote}

\subsubsection{Internet}

\textbf{Fonte:} CONATEL /\allowbreak  Speedtest Ookla (2024)

{\def\LTcaptype{none} % do not increment counter
\begin{longtable}[]{@{}ll@{}}
\toprule\noalign{}
Parâmetro & Valor \\
\midrule\noalign{}
\endhead
\bottomrule\noalign{}
\endlastfoot
Velocidade média download & 38 Mbps \\
Domicílios com internet (dept.) & 48\% \\
Tecnologia predominante & rádio \\
Opção rural & Starlink disponível (\textasciitilde USD 44/\allowbreak mês) \\
\end{longtable}
}

\subsubsection{Mercado Imobiliário e Terra
Rural}

\textbf{Fonte:} INDERT /\allowbreak  Clasificados.com.py (2024)

{\def\LTcaptype{none} % do not increment counter
\begin{longtable}[]{@{}ll@{}}
\toprule\noalign{}
Tipo & Referência \\
\midrule\noalign{}
\endhead
\bottomrule\noalign{}
\endlastfoot
Terra agrícola alta prod. (USD/\allowbreak ha) & 3,500 \\
Imóvel urbano (USD/\allowbreak m²) & 650 \\
Aluguel 2 quartos (USD/\allowbreak mês) & 260 \\
\end{longtable}
}

\begin{quote}
Valores de referência departamental. Localidades menores podem ter
preços 20--40\% abaixo da capital departamental.
\end{quote}

\subsubsection{Saúde}

\textbf{Fonte:} MSPBS /\allowbreak  IPS Paraguay (2024-2026), consolidação
departamental e proxy local

{\def\LTcaptype{none} % do not increment counter
\begin{longtable}[]{@{}ll@{}}
\toprule\noalign{}
Serviço & Disponibilidade \\
\midrule\noalign{}
\endhead
\bottomrule\noalign{}
\endlastfoot
USF /\allowbreak  Posto de Saúde & sim \\
Hospital Regional & não \\
IPS (seguro social) & sim \\
Farmácia & sim \\
Distância ao hospital de referência & \textasciitilde24 km
(Concepción) \\
\end{longtable}
}

\textbf{Principais estabelecimentos:} USF local; referencia hospitalar
em Concepción

\textbf{Observação para imigrantes:} Atendimento primario local ou em
raio curto; casos de maior complexidade seguem para Concepción.
Cobertura privada continua recomendada para especialidades e urgências
de maior complexidade.
\newpage
\secaoDiagnostico{Loma Plata}{\mapaDistrito{1605}}{Paso Barreto (Concepción) apresenta perfil operacional consolidado para
comparacao territorial, com leitura conjunta de infraestrutura, risco
hidroclimatico e capacidade de continuidade de servicos essenciais.

\subsubsection{Por que sim}

\begin{itemize}
\tightlist
\item
  Base institucional de referencia disponivel e rastreavel.
\item
  Estrutura comparativa (A-E/\allowbreak GSS) consistente para decisao relativa.
\item
  Plano de mitigacao acionavel ja definido no dossie.
\end{itemize}

\subsubsection{Por que não}

\begin{itemize}
\tightlist
\item
  Serie oficial distrital granular ainda pode apresentar lacunas
  temporais.
\item
  Sensibilidade do score exige monitoramento continuo de risco e
  conectividade.
\item
  Decisao final depende de validacao de campo e janela temporal recente.
\end{itemize}

\subsubsection{Pre-condicoes Minimas de
Decisao}

\begin{itemize}
\tightlist
\item
  Atualizar indicadores criticos em ciclo trimestral.
\item
  Confirmar trafegabilidade e contingencia hidrica/\allowbreak energetica local.
\item
  Validar checklist de resiliencia domiciliar/\allowbreak logistica antes de decisao
  final.
\end{itemize}

\subsubsection{Recomendação
Técnica}

Uso recomendado como candidato em analise multicriterio, condicionado ao
cumprimento das pre-condicoes acima. Referência atual de score: GSS 7.0;
risco predominante: moderado.

\subsubsection{}

\begin{itemize}
\tightlist
\item
  \begin{enumerate}
  \def\labelenumi{\arabic{enumi})}
  \tightlist
  \item
    Delimitacao territorial validada por georreferencia institucional.
  \end{enumerate}
\item
  \begin{enumerate}
  \def\labelenumi{\arabic{enumi})}
  \setcounter{enumi}{1}
  \tightlist
  \item
    População municipal/\allowbreak distrital referenciada por base censitaria
    nacional.
  \end{enumerate}
\item
  \begin{enumerate}
  \def\labelenumi{\arabic{enumi})}
  \setcounter{enumi}{2}
  \tightlist
  \item
    Comparabilidade intra-departamental apoiada em indicadores
    distritais oficiais.
  \end{enumerate}
\item
  \begin{enumerate}
  \def\labelenumi{\arabic{enumi})}
  \setcounter{enumi}{3}
  \tightlist
  \item
    Estado macro de conectividade viaria ancorado em informacoes de
    rodovias.
  \end{enumerate}
\item
  \begin{enumerate}
  \def\labelenumi{\arabic{enumi})}
  \setcounter{enumi}{4}
  \tightlist
  \item
    Exposicao hidrometeorologica monitorada por avisos oficiais.
  \end{enumerate}
\item
  \begin{enumerate}
  \def\labelenumi{\arabic{enumi})}
  \setcounter{enumi}{5}
  \tightlist
  \item
    Historico de contexto meteorologico apoiado em publicacoes tecnicas.
  \end{enumerate}
\item
  \begin{enumerate}
  \def\labelenumi{\arabic{enumi})}
  \setcounter{enumi}{6}
  \tightlist
  \item
    Capacidade institucional de resposta a eventos apoiada em acoes da
    SEN\footnote{Secretaría de Emergencia Nacional (Secretaria de Emergência Nacional), órgão governamental de resposta a desastres.}.
  \end{enumerate}
\item
  \begin{enumerate}
  \def\labelenumi{\arabic{enumi})}
  \setcounter{enumi}{7}
  \tightlist
  \item
    Infraestrutura energetica analisada via operador nacional.
  \end{enumerate}
\item
  \begin{enumerate}
  \def\labelenumi{\arabic{enumi})}
  \setcounter{enumi}{8}
  \tightlist
  \item
    Referências de obras e logistica nacional verificadas no MOPC.
  \end{enumerate}
\item
  \begin{enumerate}
  \def\labelenumi{\arabic{enumi})}
  \setcounter{enumi}{9}
  \tightlist
  \item
    Contexto sociopolitico institucional referenciado por autoridade
    eleitoral.
  \end{enumerate}
\item
  \begin{enumerate}
  \def\labelenumi{\arabic{enumi})}
  \setcounter{enumi}{10}
  \tightlist
  \item
    Camada cartografica de apoio eleitoral/\allowbreak territorial consultada.
  \end{enumerate}
\item
  \begin{enumerate}
  \def\labelenumi{\arabic{enumi})}
  \setcounter{enumi}{11}
  \tightlist
  \item
    Dados publicos complementares para verificacao cruzada.
  \end{enumerate}
\end{itemize}}{dist:16_Boqueron:Loma_Plata}
\begin{table}[ht!]
\small \begin{tabular}{p{3.0cm}p{0.8cm}p{6.0cm}} \toprule
\textbf{Critério} & \textbf{Nota} & \textbf{Justificativa} \\ \midrule
A - Ameaças Estratégicas & 7.2 & - \\
B - Risco Social & 7.2 & - \\
C - Riscos Naturais & 6.4 & - \\
D - Autossuficiência & 6.4 & - \\
E - Institucional & 7.6 & - \\
\midrule \textbf{GSS FINAL} & \textbf{7.0} & \textbf{Seguro (bom potencial com preparacao).} \\ \bottomrule \end{tabular}
\end{table}
\subsubsection{Vulnerabilidades e
Mitigação}\label{vuln-Loma_Plata-vulnerabilidades-e-mitigauxe7uxe3o}

\begin{itemize}
\tightlist
\item
  Exposicao hidrometeorologica moderada (45-64/\allowbreak 100).
\item
  Conectividade viaria intermediaria (50-69/\allowbreak 100).
\item
  Estoque de contingencia para 14 dias (agua/\allowbreak alimentos/\allowbreak combustivel),
  priorizado quando risco hidro=45/\allowbreak 100.
\item
  Duplicar rotas/\allowbreak logistica quando conectividade viaria=51/\allowbreak 100 for
  \textless70; manter rota secundaria mapeada e testada trimestralmente.
\item
  Plano de energia resiliente com autonomia minima de 72h
  (geracao/\allowbreak backup), revisao semestral.
\end{itemize}

Sintese aplicada: - Dados textuais consolidados com base nas fontes
oficiais acima; proximos ciclos podem aprofundar indicadores numericos
especificos por distrito.
\subsubsection{Dossiê de Campo}\label{dossie-Loma_Plata-dossiuxea-de-campo}
\begin{description}[style=nextline,leftmargin=0.5cm,font=\bfseries]
\item[Coordenadas]  23°03'S 57°07'W (geocodificado via Nominatim).
\end{description}
\textbf{Fonte climática:} NASA POWER Climatology API (período 2001-2020)

\textbf{Fonte luminosa:} estimativa world atlas

\paragraph{Irradiação Solar
(kWh/\allowbreak m²/\allowbreak dia)}\label{irradiauxe7uxe3o-solar-kwhmuxb2dia}

{\def\LTcaptype{none} % do not increment counter
\begin{longtable}[]{@{}
  >{\raggedright\arraybackslash}p{(\linewidth - 24\tabcolsep) * \real{0.0746}}
  >{\raggedright\arraybackslash}p{(\linewidth - 24\tabcolsep) * \real{0.0746}}
  >{\raggedright\arraybackslash}p{(\linewidth - 24\tabcolsep) * \real{0.0746}}
  >{\raggedright\arraybackslash}p{(\linewidth - 24\tabcolsep) * \real{0.0746}}
  >{\raggedright\arraybackslash}p{(\linewidth - 24\tabcolsep) * \real{0.0746}}
  >{\raggedright\arraybackslash}p{(\linewidth - 24\tabcolsep) * \real{0.0746}}
  >{\raggedright\arraybackslash}p{(\linewidth - 24\tabcolsep) * \real{0.0746}}
  >{\raggedright\arraybackslash}p{(\linewidth - 24\tabcolsep) * \real{0.0746}}
  >{\raggedright\arraybackslash}p{(\linewidth - 24\tabcolsep) * \real{0.0746}}
  >{\raggedright\arraybackslash}p{(\linewidth - 24\tabcolsep) * \real{0.0746}}
  >{\raggedright\arraybackslash}p{(\linewidth - 24\tabcolsep) * \real{0.0746}}
  >{\raggedright\arraybackslash}p{(\linewidth - 24\tabcolsep) * \real{0.0746}}
  >{\raggedright\arraybackslash}p{(\linewidth - 24\tabcolsep) * \real{0.1045}}@{}}
\toprule\noalign{}
\begin{minipage}[b]{\linewidth}\raggedright
Jan
\end{minipage} & \begin{minipage}[b]{\linewidth}\raggedright
Fev
\end{minipage} & \begin{minipage}[b]{\linewidth}\raggedright
Mar
\end{minipage} & \begin{minipage}[b]{\linewidth}\raggedright
Abr
\end{minipage} & \begin{minipage}[b]{\linewidth}\raggedright
Mai
\end{minipage} & \begin{minipage}[b]{\linewidth}\raggedright
Jun
\end{minipage} & \begin{minipage}[b]{\linewidth}\raggedright
Jul
\end{minipage} & \begin{minipage}[b]{\linewidth}\raggedright
Ago
\end{minipage} & \begin{minipage}[b]{\linewidth}\raggedright
Set
\end{minipage} & \begin{minipage}[b]{\linewidth}\raggedright
Out
\end{minipage} & \begin{minipage}[b]{\linewidth}\raggedright
Nov
\end{minipage} & \begin{minipage}[b]{\linewidth}\raggedright
Dez
\end{minipage} & \begin{minipage}[b]{\linewidth}\raggedright
Média
\end{minipage} \\
\midrule\noalign{}
\endhead
\bottomrule\noalign{}
\endlastfoot
6.68 & 6.07 & 5.58 & 4.63 & 3.47 & 3.00 & 3.30 & 4.09 & 4.68 & 5.49 &
6.34 & 6.61 & \textbf{5.0} \\
\end{longtable}
}

\textbf{Inclinação solar recomendada:} 23° N (anual) · 33° N (inverno
jun-ago) · 13° N (verão nov-jan)

\paragraph{Precipitação (mm/\allowbreak mês)}\label{precipitauxe7uxe3o-mmmuxeas}

{\def\LTcaptype{none} % do not increment counter
\begin{longtable}[]{@{}
  >{\raggedright\arraybackslash}p{(\linewidth - 24\tabcolsep) * \real{0.0704}}
  >{\raggedright\arraybackslash}p{(\linewidth - 24\tabcolsep) * \real{0.0704}}
  >{\raggedright\arraybackslash}p{(\linewidth - 24\tabcolsep) * \real{0.0704}}
  >{\raggedright\arraybackslash}p{(\linewidth - 24\tabcolsep) * \real{0.0704}}
  >{\raggedright\arraybackslash}p{(\linewidth - 24\tabcolsep) * \real{0.0704}}
  >{\raggedright\arraybackslash}p{(\linewidth - 24\tabcolsep) * \real{0.0704}}
  >{\raggedright\arraybackslash}p{(\linewidth - 24\tabcolsep) * \real{0.0704}}
  >{\raggedright\arraybackslash}p{(\linewidth - 24\tabcolsep) * \real{0.0704}}
  >{\raggedright\arraybackslash}p{(\linewidth - 24\tabcolsep) * \real{0.0704}}
  >{\raggedright\arraybackslash}p{(\linewidth - 24\tabcolsep) * \real{0.0704}}
  >{\raggedright\arraybackslash}p{(\linewidth - 24\tabcolsep) * \real{0.0704}}
  >{\raggedright\arraybackslash}p{(\linewidth - 24\tabcolsep) * \real{0.0704}}
  >{\raggedright\arraybackslash}p{(\linewidth - 24\tabcolsep) * \real{0.1549}}@{}}
\toprule\noalign{}
\begin{minipage}[b]{\linewidth}\raggedright
Jan
\end{minipage} & \begin{minipage}[b]{\linewidth}\raggedright
Fev
\end{minipage} & \begin{minipage}[b]{\linewidth}\raggedright
Mar
\end{minipage} & \begin{minipage}[b]{\linewidth}\raggedright
Abr
\end{minipage} & \begin{minipage}[b]{\linewidth}\raggedright
Mai
\end{minipage} & \begin{minipage}[b]{\linewidth}\raggedright
Jun
\end{minipage} & \begin{minipage}[b]{\linewidth}\raggedright
Jul
\end{minipage} & \begin{minipage}[b]{\linewidth}\raggedright
Ago
\end{minipage} & \begin{minipage}[b]{\linewidth}\raggedright
Set
\end{minipage} & \begin{minipage}[b]{\linewidth}\raggedright
Out
\end{minipage} & \begin{minipage}[b]{\linewidth}\raggedright
Nov
\end{minipage} & \begin{minipage}[b]{\linewidth}\raggedright
Dez
\end{minipage} & \begin{minipage}[b]{\linewidth}\raggedright
Total/\allowbreak ano
\end{minipage} \\
\midrule\noalign{}
\endhead
\bottomrule\noalign{}
\endlastfoot
165 & 157 & 117 & 134 & 112 & 71 & 46 & 32 & 73 & 146 & 205 & 165 &
\textbf{1428 mm} \\
\end{longtable}
}

\paragraph{Poluição Luminosa}\label{poluiuxe7uxe3o-luminosa}

{\def\LTcaptype{none} % do not increment counter
\begin{longtable}[]{@{}ll@{}}
\toprule\noalign{}
Parâmetro & Valor \\
\midrule\noalign{}
\endhead
\bottomrule\noalign{}
\endlastfoot
Escala Bortle & 3 --- Céu rural \\
Radiância artificial & 1.5 nW/\allowbreak cm²/\allowbreak sr \\
\end{longtable}
}
\paragraph{Solo (SoilGrids 2.0, média ponderada 0--30
cm)}

{\def\LTcaptype{none} % do not increment counter
\begin{longtable}[]{@{}ll@{}}
\toprule\noalign{}
Parâmetro & Valor \\
\midrule\noalign{}
\endhead
\bottomrule\noalign{}
\endlastfoot
pH (H$_2$O) & 5.75 \\
Carbono orgânico (SOC) & 15.32 g/\allowbreak kg \\
Argila & 21.42 \% \\
Areia & 54.93 \% \\
Silte (calc.) & 23.6 \% \\
Densidade aparente & 1.31 g/\allowbreak cm³ \\
\textbf{Aptidão agrícola} & \textbf{Alta} \\
\end{longtable}
}

\subsubsection{Indicadores Sociais}

\textbf{Fontes:} PNUD 2020, INE\footnote{Instituto Nacional de Estadística (Instituto Nacional de Estatística), órgão oficial de dados demográficos e censitários.} Censo 2022, INE\footnote{Instituto Nacional de Estadística (Instituto Nacional de Estatística), órgão oficial de dados demográficos e censitários.} EPHC 2023, Ministerio
Público 2024\\
\textbf{Nota:} Valores marcados como \emph{(dept.)} referem-se ao
departamento de Concepción; valores \emph{(dist.)} são específicos deste
distrito. Dados marcados \emph{(est.)} são estimativas.

{\def\LTcaptype{none} % do not increment counter
\begin{longtable}[]{@{}
  >{\raggedright\arraybackslash}p{(\linewidth - 4\tabcolsep) * \real{0.4231}}
  >{\raggedright\arraybackslash}p{(\linewidth - 4\tabcolsep) * \real{0.2692}}
  >{\raggedright\arraybackslash}p{(\linewidth - 4\tabcolsep) * \real{0.3077}}@{}}
\toprule\noalign{}
\begin{minipage}[b]{\linewidth}\raggedright
Indicador
\end{minipage} & \begin{minipage}[b]{\linewidth}\raggedright
Valor
\end{minipage} & \begin{minipage}[b]{\linewidth}\raggedright
Âmbito
\end{minipage} \\
\midrule\noalign{}
\endhead
\bottomrule\noalign{}
\endlastfoot
IDH (2020) & 0.635 & dist. (est.) \\
Ranking IDH nacional & 9/\allowbreak 18 & dept. \\
Esperança de vida & 72,4 anos & dept. \\
Escolaridade média & 7.2 anos & dist. \\
RNB per capita (USD PPA) & 8.400 & dept. \\
Pobreza monetária (\%) & 31,1\% & dept. \\
Pobreza extrema (\%) & 7,8\% & dept. \\
Índice de Gini & 0,460 & dept. \\
Acesso a água potável (\%) & 74,2\% & dept. \\
Acesso a saneamento (\%) & 71,8\% & dept. \\
Taxa de homicídios (est., /\allowbreak 100k hab) & \textasciitilde12--15
(estimativa, acima da média nacional) & dept. \\
Índice de segurança & baixo & dept. \\
População (Censo 2022) & 3.988 hab. & dist. \\
Idade mediana & 25.0 anos & dist. \\
Taxa de fecundidade & 2 & dist. \\
\end{longtable}
}
\subsubsection{Combustível}

\textbf{Referência:} PETROPAR /\allowbreak  postos locais (2024) \textbf{Tipo de
localidade:} Interior

{\def\LTcaptype{none} % do not increment counter
\begin{longtable}[]{@{}lll@{}}
\toprule\noalign{}
Combustível & USD/\allowbreak litro & Gs/\allowbreak litro (aprox.) \\
\midrule\noalign{}
\endhead
\bottomrule\noalign{}
\endlastfoot
Gasolina 93 oct & 1.01 & 7,474 \\
Gasolina 97 oct (premium) & 1.11 & 8,214 \\
Diesel & 0.94 & 6,956 \\
\end{longtable}
}

\begin{quote}
Preços podem variar ±5\% conforme posto e sazonalidade. Chaco e interior
remoto apresentam maior variação.
\end{quote}

\subsubsection{Cobertura Celular}

\textbf{Fonte:} CONATEL PY /\allowbreak  operadoras (2024)

{\def\LTcaptype{none} % do not increment counter
\begin{longtable}[]{@{}ll@{}}
\toprule\noalign{}
Parâmetro & Valor \\
\midrule\noalign{}
\endhead
\bottomrule\noalign{}
\endlastfoot
Cobertura 4G população (dept.) & 78\% \\
Cobertura 4G área rural & 55\% \\
Melhor operadora & Tigo \\
Qualidade rural & moderada \\
\end{longtable}
}

\begin{quote}
Para áreas rurais fora do núcleo urbano, recomenda-se chip Tigo como
principal e Personal como backup.
\end{quote}

\subsubsection{Internet}

\textbf{Fonte:} CONATEL /\allowbreak  Speedtest Ookla (2024)

{\def\LTcaptype{none} % do not increment counter
\begin{longtable}[]{@{}ll@{}}
\toprule\noalign{}
Parâmetro & Valor \\
\midrule\noalign{}
\endhead
\bottomrule\noalign{}
\endlastfoot
Velocidade média download & 38 Mbps \\
Domicílios com internet (dept.) & 48\% \\
Tecnologia predominante & rádio \\
Opção rural & Starlink disponível (\textasciitilde USD 44/\allowbreak mês) \\
\end{longtable}
}

\subsubsection{Mercado Imobiliário e Terra
Rural}

\textbf{Fonte:} INDERT /\allowbreak  Clasificados.com.py (2024)

{\def\LTcaptype{none} % do not increment counter
\begin{longtable}[]{@{}ll@{}}
\toprule\noalign{}
Tipo & Referência \\
\midrule\noalign{}
\endhead
\bottomrule\noalign{}
\endlastfoot
Terra agrícola alta prod. (USD/\allowbreak ha) & 3,500 \\
Imóvel urbano (USD/\allowbreak m²) & 650 \\
Aluguel 2 quartos (USD/\allowbreak mês) & 260 \\
\end{longtable}
}

\begin{quote}
Valores de referência departamental. Localidades menores podem ter
preços 20--40\% abaixo da capital departamental.
\end{quote}

\subsubsection{Saúde}

\textbf{Fonte:} MSPBS /\allowbreak  IPS Paraguay (2024-2026), consolidação
departamental e proxy local

{\def\LTcaptype{none} % do not increment counter
\begin{longtable}[]{@{}ll@{}}
\toprule\noalign{}
Serviço & Disponibilidade \\
\midrule\noalign{}
\endhead
\bottomrule\noalign{}
\endlastfoot
USF /\allowbreak  Posto de Saúde & sim \\
Hospital Regional & não \\
IPS (seguro social) & sim \\
Farmácia & sim \\
Distância ao hospital de referência & \textasciitilde65 km
(Concepción) \\
\end{longtable}
}

\textbf{Principais estabelecimentos:} USF local; referencia hospitalar
em Concepción

\textbf{Observação para imigrantes:} Atendimento primario local ou em
raio curto; casos de maior complexidade seguem para Concepción.
Cobertura privada continua recomendada para especialidades e urgências
de maior complexidade.
\newpage
\secaoDiagnostico{Mariscal Estigarribia}{\mapaConteudo{-22.0166}{-60.6000}}{Paso Horqueta (Concepción) apresenta perfil operacional consolidado para
comparacao territorial, com leitura conjunta de infraestrutura, risco
hidroclimatico e capacidade de continuidade de servicos essenciais.

\subsubsection{Por que sim}

\begin{itemize}
\tightlist
\item
  Base institucional de referencia disponivel e rastreavel.
\item
  Estrutura comparativa (A-E/\allowbreak GSS) consistente para decisao relativa.
\item
  Plano de mitigacao acionavel ja definido no dossie.
\end{itemize}

\subsubsection{Por que não}

\begin{itemize}
\tightlist
\item
  Serie oficial distrital granular ainda pode apresentar lacunas
  temporais.
\item
  Sensibilidade do score exige monitoramento continuo de risco e
  conectividade.
\item
  Decisao final depende de validacao de campo e janela temporal recente.
\end{itemize}

\subsubsection{Pre-condicoes Minimas de
Decisao}

\begin{itemize}
\tightlist
\item
  Atualizar indicadores criticos em ciclo trimestral.
\item
  Confirmar trafegabilidade e contingencia hidrica/\allowbreak energetica local.
\item
  Validar checklist de resiliencia domiciliar/\allowbreak logistica antes de decisao
  final.
\end{itemize}

\subsubsection{Recomendação
Técnica}

Uso recomendado como candidato em analise multicriterio, condicionado ao
cumprimento das pre-condicoes acima. Referência atual de score: GSS 8.0;
risco predominante: baixo.

\subsubsection{}

\begin{itemize}
\tightlist
\item
  \begin{enumerate}
  \def\labelenumi{\arabic{enumi})}
  \tightlist
  \item
    Delimitacao territorial validada por georreferencia institucional.
  \end{enumerate}
\item
  \begin{enumerate}
  \def\labelenumi{\arabic{enumi})}
  \setcounter{enumi}{1}
  \tightlist
  \item
    População municipal/\allowbreak distrital referenciada por base censitaria
    nacional.
  \end{enumerate}
\item
  \begin{enumerate}
  \def\labelenumi{\arabic{enumi})}
  \setcounter{enumi}{2}
  \tightlist
  \item
    Comparabilidade intra-departamental apoiada em indicadores
    distritais oficiais.
  \end{enumerate}
\item
  \begin{enumerate}
  \def\labelenumi{\arabic{enumi})}
  \setcounter{enumi}{3}
  \tightlist
  \item
    Estado macro de conectividade viaria ancorado em informacoes de
    rodovias.
  \end{enumerate}
\item
  \begin{enumerate}
  \def\labelenumi{\arabic{enumi})}
  \setcounter{enumi}{4}
  \tightlist
  \item
    Exposicao hidrometeorologica monitorada por avisos oficiais.
  \end{enumerate}
\item
  \begin{enumerate}
  \def\labelenumi{\arabic{enumi})}
  \setcounter{enumi}{5}
  \tightlist
  \item
    Historico de contexto meteorologico apoiado em publicacoes tecnicas.
  \end{enumerate}
\item
  \begin{enumerate}
  \def\labelenumi{\arabic{enumi})}
  \setcounter{enumi}{6}
  \tightlist
  \item
    Capacidade institucional de resposta a eventos apoiada em acoes da
    SEN\footnote{Secretaría de Emergencia Nacional (Secretaria de Emergência Nacional), órgão governamental de resposta a desastres.}.
  \end{enumerate}
\item
  \begin{enumerate}
  \def\labelenumi{\arabic{enumi})}
  \setcounter{enumi}{7}
  \tightlist
  \item
    Infraestrutura energetica analisada via operador nacional.
  \end{enumerate}
\item
  \begin{enumerate}
  \def\labelenumi{\arabic{enumi})}
  \setcounter{enumi}{8}
  \tightlist
  \item
    Referências de obras e logistica nacional verificadas no MOPC.
  \end{enumerate}
\item
  \begin{enumerate}
  \def\labelenumi{\arabic{enumi})}
  \setcounter{enumi}{9}
  \tightlist
  \item
    Contexto sociopolitico institucional referenciado por autoridade
    eleitoral.
  \end{enumerate}
\item
  \begin{enumerate}
  \def\labelenumi{\arabic{enumi})}
  \setcounter{enumi}{10}
  \tightlist
  \item
    Camada cartografica de apoio eleitoral/\allowbreak territorial consultada.
  \end{enumerate}
\item
  \begin{enumerate}
  \def\labelenumi{\arabic{enumi})}
  \setcounter{enumi}{11}
  \tightlist
  \item
    Dados publicos complementares para verificacao cruzada.
  \end{enumerate}
\end{itemize}}{dist:16_Boqueron:Mariscal_Estigarribia}
\begin{table}[ht!]
\small \begin{tabular}{p{3.0cm}p{0.8cm}p{6.0cm}} \toprule
\textbf{Critério} & \textbf{Nota} & \textbf{Justificativa} \\ \midrule
A - Ameaças Estratégicas & 8.3 & - \\
B - Risco Social & 8.7 & - \\
C - Riscos Naturais & 7.5 & - \\
D - Autossuficiência & 6.6 & - \\
E - Institucional & 8.7 & - \\
\midrule \textbf{GSS FINAL} & \textbf{8.0} & \textbf{Seguro (bom potencial com preparacao).} \\ \bottomrule \end{tabular}
\end{table}
\subsubsection{Dossiê de Campo}\label{dossie-Mariscal_Estigarribia-dossiuxea-de-campo}
\begin{description}[style=nextline,leftmargin=0.5cm,font=\bfseries]
\item[Coordenadas]  23°05'S 57°23'W (geocodificado via Nominatim).
\end{description}
\textbf{Fonte climática:} NASA POWER Climatology API (período 2001-2020)

\textbf{Fonte luminosa:} estimativa world atlas

\paragraph{Irradiação Solar
(kWh/\allowbreak m²/\allowbreak dia)}\label{irradiauxe7uxe3o-solar-kwhmuxb2dia}

{\def\LTcaptype{none} % do not increment counter
\begin{longtable}[]{@{}
  >{\raggedright\arraybackslash}p{(\linewidth - 24\tabcolsep) * \real{0.0746}}
  >{\raggedright\arraybackslash}p{(\linewidth - 24\tabcolsep) * \real{0.0746}}
  >{\raggedright\arraybackslash}p{(\linewidth - 24\tabcolsep) * \real{0.0746}}
  >{\raggedright\arraybackslash}p{(\linewidth - 24\tabcolsep) * \real{0.0746}}
  >{\raggedright\arraybackslash}p{(\linewidth - 24\tabcolsep) * \real{0.0746}}
  >{\raggedright\arraybackslash}p{(\linewidth - 24\tabcolsep) * \real{0.0746}}
  >{\raggedright\arraybackslash}p{(\linewidth - 24\tabcolsep) * \real{0.0746}}
  >{\raggedright\arraybackslash}p{(\linewidth - 24\tabcolsep) * \real{0.0746}}
  >{\raggedright\arraybackslash}p{(\linewidth - 24\tabcolsep) * \real{0.0746}}
  >{\raggedright\arraybackslash}p{(\linewidth - 24\tabcolsep) * \real{0.0746}}
  >{\raggedright\arraybackslash}p{(\linewidth - 24\tabcolsep) * \real{0.0746}}
  >{\raggedright\arraybackslash}p{(\linewidth - 24\tabcolsep) * \real{0.0746}}
  >{\raggedright\arraybackslash}p{(\linewidth - 24\tabcolsep) * \real{0.1045}}@{}}
\toprule\noalign{}
\begin{minipage}[b]{\linewidth}\raggedright
Jan
\end{minipage} & \begin{minipage}[b]{\linewidth}\raggedright
Fev
\end{minipage} & \begin{minipage}[b]{\linewidth}\raggedright
Mar
\end{minipage} & \begin{minipage}[b]{\linewidth}\raggedright
Abr
\end{minipage} & \begin{minipage}[b]{\linewidth}\raggedright
Mai
\end{minipage} & \begin{minipage}[b]{\linewidth}\raggedright
Jun
\end{minipage} & \begin{minipage}[b]{\linewidth}\raggedright
Jul
\end{minipage} & \begin{minipage}[b]{\linewidth}\raggedright
Ago
\end{minipage} & \begin{minipage}[b]{\linewidth}\raggedright
Set
\end{minipage} & \begin{minipage}[b]{\linewidth}\raggedright
Out
\end{minipage} & \begin{minipage}[b]{\linewidth}\raggedright
Nov
\end{minipage} & \begin{minipage}[b]{\linewidth}\raggedright
Dez
\end{minipage} & \begin{minipage}[b]{\linewidth}\raggedright
Média
\end{minipage} \\
\midrule\noalign{}
\endhead
\bottomrule\noalign{}
\endlastfoot
6.68 & 6.07 & 5.58 & 4.63 & 3.47 & 3.00 & 3.30 & 4.09 & 4.68 & 5.49 &
6.34 & 6.61 & \textbf{5.0} \\
\end{longtable}
}

\textbf{Inclinação solar recomendada:} 23° N (anual) · 33° N (inverno
jun-ago) · 13° N (verão nov-jan)

\paragraph{Precipitação (mm/\allowbreak mês)}\label{precipitauxe7uxe3o-mmmuxeas}

{\def\LTcaptype{none} % do not increment counter
\begin{longtable}[]{@{}
  >{\raggedright\arraybackslash}p{(\linewidth - 24\tabcolsep) * \real{0.0704}}
  >{\raggedright\arraybackslash}p{(\linewidth - 24\tabcolsep) * \real{0.0704}}
  >{\raggedright\arraybackslash}p{(\linewidth - 24\tabcolsep) * \real{0.0704}}
  >{\raggedright\arraybackslash}p{(\linewidth - 24\tabcolsep) * \real{0.0704}}
  >{\raggedright\arraybackslash}p{(\linewidth - 24\tabcolsep) * \real{0.0704}}
  >{\raggedright\arraybackslash}p{(\linewidth - 24\tabcolsep) * \real{0.0704}}
  >{\raggedright\arraybackslash}p{(\linewidth - 24\tabcolsep) * \real{0.0704}}
  >{\raggedright\arraybackslash}p{(\linewidth - 24\tabcolsep) * \real{0.0704}}
  >{\raggedright\arraybackslash}p{(\linewidth - 24\tabcolsep) * \real{0.0704}}
  >{\raggedright\arraybackslash}p{(\linewidth - 24\tabcolsep) * \real{0.0704}}
  >{\raggedright\arraybackslash}p{(\linewidth - 24\tabcolsep) * \real{0.0704}}
  >{\raggedright\arraybackslash}p{(\linewidth - 24\tabcolsep) * \real{0.0704}}
  >{\raggedright\arraybackslash}p{(\linewidth - 24\tabcolsep) * \real{0.1549}}@{}}
\toprule\noalign{}
\begin{minipage}[b]{\linewidth}\raggedright
Jan
\end{minipage} & \begin{minipage}[b]{\linewidth}\raggedright
Fev
\end{minipage} & \begin{minipage}[b]{\linewidth}\raggedright
Mar
\end{minipage} & \begin{minipage}[b]{\linewidth}\raggedright
Abr
\end{minipage} & \begin{minipage}[b]{\linewidth}\raggedright
Mai
\end{minipage} & \begin{minipage}[b]{\linewidth}\raggedright
Jun
\end{minipage} & \begin{minipage}[b]{\linewidth}\raggedright
Jul
\end{minipage} & \begin{minipage}[b]{\linewidth}\raggedright
Ago
\end{minipage} & \begin{minipage}[b]{\linewidth}\raggedright
Set
\end{minipage} & \begin{minipage}[b]{\linewidth}\raggedright
Out
\end{minipage} & \begin{minipage}[b]{\linewidth}\raggedright
Nov
\end{minipage} & \begin{minipage}[b]{\linewidth}\raggedright
Dez
\end{minipage} & \begin{minipage}[b]{\linewidth}\raggedright
Total/\allowbreak ano
\end{minipage} \\
\midrule\noalign{}
\endhead
\bottomrule\noalign{}
\endlastfoot
164 & 151 & 121 & 135 & 108 & 67 & 43 & 29 & 68 & 141 & 209 & 165 &
\textbf{1405 mm} \\
\end{longtable}
}

\paragraph{Poluição Luminosa}\label{poluiuxe7uxe3o-luminosa}

{\def\LTcaptype{none} % do not increment counter
\begin{longtable}[]{@{}ll@{}}
\toprule\noalign{}
Parâmetro & Valor \\
\midrule\noalign{}
\endhead
\bottomrule\noalign{}
\endlastfoot
Escala Bortle & 3 --- Céu rural \\
Radiância artificial & 1.5 nW/\allowbreak cm²/\allowbreak sr \\
\end{longtable}
}
\paragraph{Solo (SoilGrids 2.0, média ponderada 0--30
cm)}

{\def\LTcaptype{none} % do not increment counter
\begin{longtable}[]{@{}ll@{}}
\toprule\noalign{}
Parâmetro & Valor \\
\midrule\noalign{}
\endhead
\bottomrule\noalign{}
\endlastfoot
pH (H$_2$O) & 5.88 \\
Carbono orgânico (SOC) & 16.02 g/\allowbreak kg \\
Argila & 19.35 \% \\
Areia & 53.9 \% \\
Silte (calc.) & 26.8 \% \\
Densidade aparente & 1.33 g/\allowbreak cm³ \\
\textbf{Aptidão agrícola} & \textbf{Alta} \\
\end{longtable}
}

\subsubsection{Indicadores Sociais}

\textbf{Fontes:} PNUD 2020, INE\footnote{Instituto Nacional de Estadística (Instituto Nacional de Estatística), órgão oficial de dados demográficos e censitários.} Censo 2022, INE\footnote{Instituto Nacional de Estadística (Instituto Nacional de Estatística), órgão oficial de dados demográficos e censitários.} EPHC 2023, Ministerio
Público 2024\\
\textbf{Nota:} Valores marcados como \emph{(dept.)} referem-se ao
departamento de Concepción; valores \emph{(dist.)} são específicos deste
distrito. Dados marcados \emph{(est.)} são estimativas.

{\def\LTcaptype{none} % do not increment counter
\begin{longtable}[]{@{}
  >{\raggedright\arraybackslash}p{(\linewidth - 4\tabcolsep) * \real{0.4231}}
  >{\raggedright\arraybackslash}p{(\linewidth - 4\tabcolsep) * \real{0.2692}}
  >{\raggedright\arraybackslash}p{(\linewidth - 4\tabcolsep) * \real{0.3077}}@{}}
\toprule\noalign{}
\begin{minipage}[b]{\linewidth}\raggedright
Indicador
\end{minipage} & \begin{minipage}[b]{\linewidth}\raggedright
Valor
\end{minipage} & \begin{minipage}[b]{\linewidth}\raggedright
Âmbito
\end{minipage} \\
\midrule\noalign{}
\endhead
\bottomrule\noalign{}
\endlastfoot
IDH (2020) & 0.640 & dist. (est.) \\
Ranking IDH nacional & 9/\allowbreak 18 & dept. \\
Esperança de vida & 72,4 anos & dept. \\
Escolaridade média & 7.4 anos & dist. \\
RNB per capita (USD PPA) & 8.400 & dept. \\
Pobreza monetária (\%) & 31,1\% & dept. \\
Pobreza extrema (\%) & 7,8\% & dept. \\
Índice de Gini & 0,460 & dept. \\
Acesso a água potável (\%) & 74,2\% & dept. \\
Acesso a saneamento (\%) & 71,8\% & dept. \\
Taxa de homicídios (est., /\allowbreak 100k hab) & \textasciitilde12--15
(estimativa, acima da média nacional) & dept. \\
Índice de segurança & baixo & dept. \\
População (Censo 2022) & 7.646 hab. & dist. \\
Idade mediana & 27.0 anos & dist. \\
Taxa de fecundidade & 2 & dist. \\
\end{longtable}
}
\subsubsection{Combustível}

\textbf{Referência:} PETROPAR /\allowbreak  postos locais (2024) \textbf{Tipo de
localidade:} Interior

{\def\LTcaptype{none} % do not increment counter
\begin{longtable}[]{@{}lll@{}}
\toprule\noalign{}
Combustível & USD/\allowbreak litro & Gs/\allowbreak litro (aprox.) \\
\midrule\noalign{}
\endhead
\bottomrule\noalign{}
\endlastfoot
Gasolina 93 oct & 1.01 & 7,474 \\
Gasolina 97 oct (premium) & 1.11 & 8,214 \\
Diesel & 0.94 & 6,956 \\
\end{longtable}
}

\begin{quote}
Preços podem variar ±5\% conforme posto e sazonalidade. Chaco e interior
remoto apresentam maior variação.
\end{quote}

\subsubsection{Cobertura Celular}

\textbf{Fonte:} CONATEL PY /\allowbreak  operadoras (2024)

{\def\LTcaptype{none} % do not increment counter
\begin{longtable}[]{@{}ll@{}}
\toprule\noalign{}
Parâmetro & Valor \\
\midrule\noalign{}
\endhead
\bottomrule\noalign{}
\endlastfoot
Cobertura 4G população (dept.) & 78\% \\
Cobertura 4G área rural & 55\% \\
Melhor operadora & Tigo \\
Qualidade rural & moderada \\
\end{longtable}
}

\begin{quote}
Para áreas rurais fora do núcleo urbano, recomenda-se chip Tigo como
principal e Personal como backup.
\end{quote}

\subsubsection{Internet}

\textbf{Fonte:} CONATEL /\allowbreak  Speedtest Ookla (2024)

{\def\LTcaptype{none} % do not increment counter
\begin{longtable}[]{@{}ll@{}}
\toprule\noalign{}
Parâmetro & Valor \\
\midrule\noalign{}
\endhead
\bottomrule\noalign{}
\endlastfoot
Velocidade média download & 38 Mbps \\
Domicílios com internet (dept.) & 48\% \\
Tecnologia predominante & rádio \\
Opção rural & Starlink disponível (\textasciitilde USD 44/\allowbreak mês) \\
\end{longtable}
}

\subsubsection{Mercado Imobiliário e Terra
Rural}

\textbf{Fonte:} INDERT /\allowbreak  Clasificados.com.py (2024)

{\def\LTcaptype{none} % do not increment counter
\begin{longtable}[]{@{}ll@{}}
\toprule\noalign{}
Tipo & Referência \\
\midrule\noalign{}
\endhead
\bottomrule\noalign{}
\endlastfoot
Terra agrícola alta prod. (USD/\allowbreak ha) & 3,500 \\
Imóvel urbano (USD/\allowbreak m²) & 650 \\
Aluguel 2 quartos (USD/\allowbreak mês) & 260 \\
\end{longtable}
}

\begin{quote}
Valores de referência departamental. Localidades menores podem ter
preços 20--40\% abaixo da capital departamental.
\end{quote}

\subsubsection{Saúde}

\textbf{Fonte:} MSPBS /\allowbreak  IPS Paraguay (2024-2026), consolidação
departamental e proxy local

{\def\LTcaptype{none} % do not increment counter
\begin{longtable}[]{@{}ll@{}}
\toprule\noalign{}
Serviço & Disponibilidade \\
\midrule\noalign{}
\endhead
\bottomrule\noalign{}
\endlastfoot
USF /\allowbreak  Posto de Saúde & sim \\
Hospital Regional & não \\
IPS (seguro social) & sim \\
Farmácia & sim \\
Distância ao hospital de referência & \textasciitilde47 km
(Concepción) \\
\end{longtable}
}

\textbf{Principais estabelecimentos:} USF local; referencia hospitalar
em Concepción

\textbf{Observação para imigrantes:} Atendimento primario local ou em
raio curto; casos de maior complexidade seguem para Concepción.
Cobertura privada continua recomendada para especialidades e urgências
de maior complexidade.
